\documentclass[11pt,a4paper]{article}

\usepackage[a4paper,margin=2.3cm]{geometry}
\usepackage{fontspec}
\usepackage{titlesec}
\usepackage{longtable}
\usepackage{booktabs}
\usepackage{array}
\usepackage{xcolor}
\usepackage[most]{tcolorbox}
\usepackage{hyperref}
\usepackage{enumitem}
\usepackage{fancyhdr}
\usepackage{amssymb}

\setmainfont{Sylfaen}

\definecolor{brandgreen}{HTML}{1F6F5C}
\definecolor{quoteblue}{HTML}{1D4E89}
\definecolor{termorange}{HTML}{B5651D}
\definecolor{importred}{HTML}{B03A2E}

\hypersetup{
  colorlinks=true,
  linkcolor=brandgreen,
  urlcolor=brandgreen,
  pdftitle={Lowe et al. 2021 (Brazil) -- ქართული თარგმანი და შენიშვნები},
  pdfauthor={სერგი კონიაშვილი}
}

\titleformat{\section}{\large\bfseries\color{brandgreen}}{\thesection}{0.5em}{}
\titleformat{\subsection}{\normalsize\bfseries}{\thesubsection}{0.5em}{}

\setlength{\headheight}{14pt}
\pagestyle{fancy}
\fancyhf{}
\fancyhead[L]{Lowe et al. (2021), Lancet Planetary Health}
\fancyhead[R]{ქართული თარგმანი სასწავლად}
\fancyfoot[C]{\thepage}
\setlength{\emergencystretch}{3em}
\setlength{\parskip}{0.4em}

\newtcolorbox{etl}{enhanced,breakable,colback=brandgreen!6,colframe=brandgreen,
  fonttitle=\bfseries,title={ETL: მნიშვნელოვანია ჩემი სისტემისთვის}}
\newtcolorbox{citat}{enhanced,breakable,colback=quoteblue!5,colframe=quoteblue,
  fonttitle=\bfseries,title={ციტატა თეზისისთვის (ორიგინალი ინგლისურად)}}
\newtcolorbox{termin}{enhanced,breakable,colback=termorange!7,colframe=termorange,
  fonttitle=\bfseries,title={ტერმინები}}
\newtcolorbox{skipbox}{enhanced,breakable,colback=gray!7,colframe=gray!55,
  fonttitle=\bfseries,title={ნაკლებად მნიშვნელოვანია შენთვის: საკმარისია ზოგადი გაგება}}
\newtcolorbox{alarm}{enhanced,breakable,colback=importred!5,colframe=importred,
  fonttitle=\bfseries,title={ყურადღება: ეს შენს სისტემას პირდაპირ ეხება}}
\newcommand{\imp}[1]{\textbf{\color{importred}#1}}
\newcommand{\pg}[1]{{\small\color{gray}[ორიგინალი, გვ.~#1]}}

\begin{document}

\begin{center}
  {\Large\bfseries Combined effects of hydrometeorological hazards and urbanisation on dengue risk in Brazil: a spatiotemporal modelling study}\\[0.6em]
  {\large ჰიდრომეტეოროლოგიური საფრთხეებისა და ურბანიზაციის ერთობლივი გავლენა დენგეს რისკზე ბრაზილიაში: სივრცულ-დროითი მოდელირების კვლევა}\\[0.8em]
  Lowe R., Lee S.A., O'Reilly K.M., Brady O.J., Bastos L., Carrasco-Escobar G., de Castro Catão R.,
  Colón-González F.J., Barcellos C., Carvalho M.S., Blangiardo M., Rue H., Gasparrini A.\\
  \textit{The Lancet Planetary Health} 5(4): e209--e219 (2021). ლიცენზია: CC BY 4.0\\
  DOI: \href{https://doi.org/10.1016/S2542-5196(20)30292-8}{10.1016/S2542-5196(20)30292-8}\\
  მონაცემები და კოდი: \href{https://github.com/drrachellowe/hydromet_dengue}{github.com/drrachellowe/hydromet\_dengue},
  Zenodo: \href{https://doi.org/10.5281/zenodo.4632205}{10.5281/zenodo.4632205}
\end{center}

\noindent\rule{\linewidth}{0.4pt}

\noindent\textbf{ეს დოკუმენტი სასწავლო თარგმანია და ოფიციალური არ არის.} სტატია ღიაა (CC BY 4.0), ამიტომ წყაროს მითითებით
თარგმნა დასაშვებია. თეზისში \imp{ყოველთვის ორიგინალი ინგლისური ტექსტი} ციტირე. გვერდები ჟურნალის ნუმერაციითაა მითითებული
(e209--e219), როგორც ციტირებისას გჭირდება. სრულადაა თარგმნილი: Summary, Research in context, \textbf{Methods (ყველაზე მნიშვნელოვანი)},
Limitations, Data sharing. Introduction, Results და Discussion შემოკლებულია, ეს ადგილები \textit{[შემოკლებით]} არის მონიშნული.

\subsection*{რატომ არის ეს სტატია შენთვის მნიშვნელოვანი}
ეს \textbf{მოდელირების სტატიაა}, ისევე როგორც Suttirat 2026, ოღონდ სხვა ქვეყანაზე (ბრაზილია) და სხვა ავტორების მიერ.
შენთვის მისი მთავარი ღირებულება ისაა, რომ აჩვენებს: Suttirat-ის მონაცემთა მომზადების ნაბიჯები \imp{ტიპურია} და არა ერთი
სტატიის თავისებურება (არგუმენტაციის ჯაჭვის მე-3 ნაბიჯი). კონკრეტულად:
\begin{enumerate}[leftmargin=1.5em]
  \item \textbf{იგივე ნაბიჯები}: შემთხვევები სამინისტროდან + მოსახლეობა $\rightarrow$ ინციდენტობა 100\,000-ზე; გრიდული კლიმატი
    $\rightarrow$ რეგიონზე გასაშუალოება; \textbf{Tmin/Tmax}; \textbf{lag-ები} (აქ 0--6 თვე); სოციალურ-ეკონომიკური ცვლადები.
  \item \textbf{ხუთი სხვადასხვა წყარო}, ხუთი სხვადასხვა ფორმატითა და დროითი რეზოლუციით: სწორედ ეს პრობლემაა, რომელსაც შენი სისტემა წყვეტს.
  \item \textbf{ახალი მოთხოვნები} შენი სისტემისთვის: კლიმატის მონაცემები \imp{საწყის თარიღამდეც} (lag-ისთვის), ფართობით შეწონილი
    გასაშუალოება, გვალვის ინდექსი.
\end{enumerate}

\subsection*{ფერების ლეგენდა}
\textcolor{brandgreen}{\textbf{მწვანე}}: მნიშვნელოვანია შენი სისტემისთვის;
\textcolor{quoteblue}{\textbf{ლურჯი}}: ციტატა თეზისისთვის;
\textcolor{termorange}{\textbf{ნარინჯისფერი}}: ტერმინები;
\textcolor{importred}{\textbf{წითელი ყუთი}}: შენი სისტემის პრობლემა ან შესამოწმებელი საკითხი;
\textcolor{gray}{\textbf{ნაცრისფერი}}: ნაკლებად მნიშვნელოვანი.

%==========================================================
\section{Summary (რეზიუმე)} \pg{e209}

\textbf{ფონი.} ცნობილია, რომ ტემპერატურისა და ნალექის მახასიათებლები დენგეს გადაცემის სეზონურობაზე მოქმედებს. თუმცა ძლიერი
გვალვისა და უკიდურესად ნოტიო პირობების გავლენა დენგეს ეპიდემიების დროსა და ინტენსივობაზე ცუდადაა შესწავლილი. ამ კვლევაში
ჩვენი მიზანი იყო, სივრცულ-დროითი მოდელის საშუალებით რაოდენობრივად შეგვეფასებინა ექსტრემალური ჰიდრომეტეოროლოგიური
საფრთხეების არაწრფივი და დაგვიანებული გავლენა დენგეს რისკზე ბრაზილიაში, ურბანიზაციის დონის მიხედვით.

\textbf{მეთოდები.} ჩვენ გავაერთიანეთ განაწილებული დაგვიანების არაწრფივი მოდელები (DLNM) სივრცულ-დროით ბაიესურ იერარქიულ
მოდელთან, რათა განგვესაზღვრა ზემოქმედება-დაგვიანება-პასუხის კავშირი დენგეს ფარდობით რისკსა (RR) და გვალვის სიმძიმის ინდექსს
შორის. მოდელი მოვარგეთ ბრაზილიის \textbf{558 მიკრორეგიონის თვიურ} დენგეს შემთხვევებს 2001 წლის იანვრიდან 2019 წლის იანვრამდე%
\footnote{შენიშვნა: რეზიუმეში წერია „January, 2019“, Methods-ში კი „December, 2019“ (გვ.~e211). სტატიაში ეს შეუსაბამობაა;
თეზისში Methods-ის თარიღი გამოიყენე: 2001-01 -- 2019-12.},
გავითვალისწინეთ დაუკვირვებელი შემრევი ფაქტორები, სივრცული ავტოკორელაცია, სეზონურობა და წლებს შორის ცვალებადობა. RR-ის
ცვლილება ურბანიზაციის დონის მიხედვით შევაფასეთ გვალვის ინდექსსა და ურბანიზაციას შორის ურთიერთქმედებით. ასევე შევაფასეთ
ჰიდრომეტეოროლოგიური საფრთხეების გავლენა იმ რაიონებში, სადაც წყალმომარაგების შეფერხებები ხშირია.

\textbf{შედეგები.} მონაცემთა ნაკრები მოიცავდა 12\,895\,293 დენგეს შემთხვევას (2001--2019). დენგეს რისკი გაიზარდა უკიდურესად
ნოტიო პირობებიდან \textbf{0--3 თვის} შემდეგ (მაქსიმალური RR 1 თვის lag-ზე: 1,56 [95\% CI 1,41--1,73]) და გვალვიდან
\textbf{3--5 თვის} შემდეგ (მაქსიმალური RR 4 თვის lag-ზე: 1,43 [1,22--1,67]). ურბანიზაციასთან ურთიერთქმედებამ მოდელი გააუმჯობესა:
ნოტიო პირობებში რისკი უფრო მაღალი იყო სოფლის ტიპის რაიონებში, ექსტრემალური გვალვის შემდეგ კი, პირიქით, მაღალურბანიზებულ
რაიონებში. ექსტრემალური გვალვის შემდეგ რისკი მაღალი იყო იქ, სადაც წყალმომარაგება ხშირად წყდებოდა.

\textbf{ინტერპრეტაცია.} ნოტიო პირობებსაც და ექსტრემალურ გვალვასაც შეუძლია დენგეს რისკის გაზრდა, ოღონდ სხვადასხვა დაგვიანებით.
გვალვის რისკი განსაკუთრებით მაღალია ცუდად მომსახურე ქალაქებში, სადაც გვალვისას წყალი წყდება და ხალხი მას იმპროვიზირებულ
ჭურჭელში ინახავს. კოღოების კონტროლი მხოლოდ ნოტიო და თბილ სეზონში კი არა, გვალვის პერიოდშიც საჭიროა.

\begin{termin}
\textbf{ჰიდრომეტეოროლოგიური საფრთხე}: წყალთან დაკავშირებული ექსტრემალური ამინდი (გვალვა, ძლიერი ნალექი, წყალდიდობა).\\
\textbf{RR (Relative Risk, ფარდობითი რისკი)}: რამდენჯერ მეტია რისკი ნორმალურ პირობებთან შედარებით. RR = 1,56 ნიშნავს 56\%-ით მეტს.\\
\textbf{DLNM (Distributed Lag Non-linear Model)}: მოდელი, რომელიც ერთდროულად აფასებს, \textit{რამდენად} (არაწრფივად) და
\textit{რა დაგვიანებით} (lag) მოქმედებს ცვლადი შედეგზე.\\
\textbf{მიკრორეგიონი}: ბრაზილიის ადმინისტრაციული ერთეული, მუნიციპალიტეტების ჯგუფი დიდი ქალაქის გარშემო (558 ცალი).
\end{termin}

%==========================================================
\section{Introduction (შესავალი) \textit{[შემოკლებით]}} \pg{e209--e211}

დენგე WHO-ს მიერ გლობალური ჯანმრთელობის ათ უმთავრეს საფრთხეს შორისაა დასახელებული. ბრაზილიაში 2019 წელს 1,5 მილიონზე მეტი
შემთხვევა დარეგისტრირდა. გადაცემა სამხრეთით, მაღალმთიან ქალაქებსა და ამაზონიაში ვრცელდება გარემოს ცვლილების, ურბანიზაციისა
და რეგიონებს შორის კავშირების გაუმჯობესების გამო.

ტემპერატურა მოქმედებს კოღოს განვითარებაზე, გამრავლებაზე, გადარჩენაზე, კბენის სიხშირესა და ვირუსის რეპლიკაციაზე (ოპტიმუმი
26--29$^\circ$C). ნალექის გავლენა უფრო რთულია: წვიმა ქმნის ლარვების საბინადროს წყლით სავსე ჭურჭელში, მაგრამ ჭარბი წვიმა
ლარვებს რეცხავს. გვალვისას ხალხი წყალს ინახავს, რაც დამატებით საბინადროს ქმნის. ბოლო ათწლეულში ბრაზილიაში რამდენიმე
ექსტრემალური გვალვა და წყალდიდობა მოხდა (ჩრდილო-აღმოსავლეთი 2010--16, სან-პაულო 2014--15).

\begin{citat}
„The effect of climate variability and climate change on dengue transmission is complex, non-linear, and often delayed by several
weeks to months, which limits the inferences that can be made from traditional linear modelling methods.“ \pg{e211}
\end{citat}

\begin{etl}
ეს ციტატა \textbf{მეორე დამოუკიდებელი წყაროა} (Suttirat-ის შემდეგ) იმისა, რომ კავშირები \textbf{არაწრფივი და დაგვიანებულია}.
შენთვის ორ რამეს ნიშნავს:
\begin{enumerate}[leftmargin=1.5em]
  \item \textbf{lag სვეტები} (Gemini დავალება B) მეცნიერულად დასაბუთებულია: ორივე სტატია lag-ებს იყენებს.
  \item შენი სისტემის Pearson/Spearman კორელაცია მხოლოდ \imp{საორიენტაციოა} (exploratory). თეზისში ეს შეზღუდვა ამ ციტატით დაასაბუთე:
    სისტემა მოდელს არ ცვლის, ის მოდელისთვის მონაცემებს ამზადებს.
\end{enumerate}
\end{etl}

%==========================================================
\section{Research in context (კვლევის კონტექსტი)} \pg{e210}

\textbf{მტკიცებულებები ამ კვლევამდე.} ცნობილია, რომ ტემპერატურა და ნალექი დენგეს გადაცემის სიდიდესა და სეზონურობაზე
მოქმედებს. თუმცა ცუდადაა შესწავლილი, რა დროის განმავლობაში შეუძლია ექსტრემალურ მოვლენებს (გვალვა, უკიდურესი სინოტივე)
ეპიდემიის დროსა და ინტენსივობაზე გავლენა. მეტეოროლოგებმა გვალვის რამდენიმე ინდექსი შექმნეს, მაგრამ მათი გამოყენება დენგეს
კონტროლში იშვიათია. PubMed-ში ძებნამ (2020-10-15; „dengue“, „drought“, „model“) მხოლოდ ერთი კვლევა იპოვა (ბარბადოსი).

\textbf{ამ კვლევის დამატებითი ღირებულება.} ექსტრემალური გვალვა დაკავშირებულია დენგეს რისკის ზრდასთან 3--5 თვის შემდეგ,
უკიდურესად ნოტიო პირობები კი 0--3 თვის შემდეგ. ეს ადასტურებს ბარბადოსის შედეგებს კლიმატისა და ურბანიზაციის ფართო დიაპაზონზე.

\textbf{დასკვნები.} კოღოების კონტროლი საჭიროა იმ ქალაქებშიც, სადაც წყალმომარაგება წყდება, და არა მხოლოდ წვიმიან სეზონში;
გრძელვადიან პერსპექტივაში მთავრობებმა სტაბილურ წყალმომარაგებაში უნდა ჩადონ ინვესტიცია.

\begin{skipbox}
ეს ნაწილი ეპიდემიოლოგიურ დასკვნებს ეხება. შენთვის საკმარისია იცოდე, რომ \textbf{lag 0--6 თვე} და \textbf{ექსტრემალური მოვლენები}
მნიშვნელოვანია. Introduction-ში ავტორები ასევე წერენ, რომ შედეგები \textbf{ადრეული გაფრთხილების სისტემებისთვისაა} (early
warning systems) მნიშვნელოვანი, ანუ მოდელს რეგულარულად ახალი მონაცემები სჭირდება (გვ.~e211).
\end{skipbox}

%==========================================================
\section{Methods (მეთოდები): სრული თარგმანი} \pg{e211--e214}

\begin{etl}
\textbf{ეს ყველაზე მნიშვნელოვანი ნაწილია.} აქ ჩანს ყველა წყარო და ყველა მონაცემთა მომზადების ნაბიჯი. თითოეული ქვეთავის შემდეგ
შენიშნულია, რას აკეთებს ამ ნაბიჯიდან შენი სისტემა. სრული შედარება ბოლოს, ცხრილშია (ნაწილი~\ref{sec:shedareba}).
\end{etl}

\subsection*{კვლევის არეალი და დენგეს მონაცემები} \pg{e211}
ბრაზილია მსოფლიოში მოსახლეობით მეექვსე ქვეყანაა (209 მილიონზე მეტი). მას აქვს ხუთი გეოპოლიტიკური რეგიონი, 27 შტატი და 5570
მუნიციპალიტეტი, რომლებიც \textbf{558 მიკრორეგიონადაა} დაჯგუფებული. ჩვენ მივიღეთ \textbf{თვიური} დარეგისტრირებული დენგეს
შემთხვევები 558 მიკრორეგიონისთვის 2001 წლის იანვრიდან 2019 წლის დეკემბრამდე \textbf{სავალდებულოდ შეტყობინებადი დაავადებების
საინფორმაციო სისტემიდან} (SINAN), რომელიც ჯანდაცვის სამინისტროს საინფორმაციო დეპარტამენტის (\textbf{DATASUS}) მეშვეობით
თავისუფლად ხელმისაწვდომია. ბრაზილიის ჯანდაცვის სამინისტრო დენგეს თვიურ ინციდენტობას განსაზღვრავს როგორც
\textbf{ახალი შემთხვევების რაოდენობას 100\,000 მცხოვრებზე თვეში}. თითოეული მიკრორეგიონის ინციდენტობის გამოსათვლელად
მივიღეთ \textbf{წლიური მოსახლეობის} შეფასებები 2001--2019 წლებისთვის ბრაზილიის გეოგრაფიისა და სტატისტიკის ინსტიტუტიდან
(\textbf{IBGE}), DATASUS-ის მეშვეობით.

\begin{etl}
\textbf{შენს სისტემასთან შედარება:} \textbf{თვიური შემთხვევები + წლიური მოსახლეობა $\rightarrow$ ინციდენტობა 100\,000-ზე} ზუსტად
შენი FR-17-ია. ეს უკვე \textbf{მეორე სტატიაა} (Suttirat-ის შემდეგ), რომელიც ამ ნაბიჯს აკეთებს, ანუ ნაბიჯი ტიპურია.
\textbf{განსხვავება}: აქ ერთეული მიკრორეგიონია (558), შენთან ქვეყანა. ავტორები არ წერენ, წლიური მოსახლეობა თვეებზე როგორ
გადაიყვანეს (ინტერპოლაცია თუ არა), ამიტომ ამაზე არაფერი მოიგონო.
\end{etl}

\subsection*{მეტეოროლოგიური მონაცემები} \pg{e211}
ამ კვლევაში გამოვიყენეთ \textbf{პალმერის გვალვის სიმძიმის ინდექსი} (PDSI), რომელიც გვალვისა და არიდულობის გრძელვადიანი
ცვლილებების მონიტორინგის ყველაზე გავრცელებული სტანდარტიზებული ინდექსია. \textbf{დღიური მინიმალური ტემპერატურის თვიური საშუალო}
(Tmin, $^\circ$C), \textbf{მაქსიმალური ტემპერატურა} (Tmax, $^\circ$C) და თვითკალიბრებადი PDSI მივიღეთ
\textbf{Climatic Research Unit gridded Time Series}-დან (CRU TS, ვერსია 4.04), 2000 წლის იანვრიდან 2019 წლის დეკემბრამდე,
\textbf{0,5$^\circ$ $\times$ 0,5$^\circ$} სივრცული რეზოლუციით \imp{(2000 წლის მონაცემები ამოვიღეთ იმისთვის, რომ დენგეს პირველ
დაკვირვებამდე, 2001 წლის იანვრამდე, lag-ის პერიოდი გვქონოდა)}. გრიდული მონაცემები თითოეულ მიკრორეგიონზე
\textbf{აგრეგირდა R-ის პაკეტით exactextractr} (R ვერსია 4.0.2): მიკრორეგიონის შიგნით მდებარე გრიდის უჯრედების საშუალოს
გამოთვლით. \textbf{ნაწილობრივ დაფარული უჯრედები შეიწონა} იმ პროცენტით, რომელიც მიკრორეგიონის შიგნით ხვდებოდა.

PDSI მეტეოროლოგიური გვალვის ერთ-ერთი ყველაზე გავრცელებული საზომია: ის აჩვენებს რეგიონის სიმშრალეს ნორმალურ პირობებთან
შედარებით. PDSI ითვლება ნიადაგის ტენიანობით, მოსალოდნელი აორთქლებით (ტემპერატურისა და დღის ხანგრძლივობის მიხედვით) და
ნალექით. თვითკალიბრებადი PDSI თითოეული ადგილისთვის საკუთარ „ნორმას“ ითვლის, ამიტომ ადგილებს შორის შედარებადია. ინდექსის
დიაპაზონია \textbf{--10 (ძალიან მშრალი)}-დან \textbf{10 (ძალიან ნოტიო)}-მდე; \textbf{--4-ზე ნაკლები ან 4-ზე მეტი} მნიშვნელობა
ექსტრემალურად ითვლება. ბრაზილიაში 2010 წლიდან ექსტრემალური და ხანგრძლივი გვალვა იყო რამდენიმე შტატში. მინიმალური ტემპერატურა
ქვეყნის მასშტაბით ძლიერ განსხვავდება: ტროპიკულ ჩრდილოეთში მუდმივად მაღალია, ზომიერ სამხრეთში ზამთარი ცივია.

\begin{citat}
„The monthly mean daily minimum temperature (Tmin; °C), maximum temperature (Tmax; °C), and the self-calibrated PDSI were obtained
from the Climatic Research Unit gridded Time Series (version 4.04), for the period January, 2000, to December, 2019, at a spatial
resolution of 0·5° longitude × 0·5° latitude (data for 2000 were extracted to allow for a lag period before the first dengue
observation in January, 2001). The gridded datasets were aggregated to each microregion using the exactextractr package in R
(version 4.0.2), by calculating the mean of grid boxes lying within each microregion. Grid boxes that were partially covered by
the microregion were weighted by the percentage that lay within the microregion.“ \pg{e211}
\end{citat}

\begin{alarm}
\textbf{სამი რამ, რაც შენს სისტემას პირდაპირ ეხება:}
\begin{enumerate}[leftmargin=1.5em]
  \item \imp{კლიმატის მონაცემები საწყის თარიღამდე (lag buffer).} ავტორებმა კლიმატი \textbf{ერთი წლით ადრე} ჩამოტვირთეს, რომ
    პირველ თვეებსაც ჰქონოდათ lag-ის მნიშვნელობები. თუ შენი სისტემა lag სვეტებს (დავალება B) მხოლოდ მოთხოვნილი პერიოდის
    მონაცემებით ითვლის, პირველი $k$ თვის \texttt{\_lag$k$} სვეტი \textbf{ცარიელი} იქნება. გამოსავალი: კლიმატი ჩამოიტვირთოს
    \texttt{start\_date -- max(lag\_periods)} თვიდან. ეს Gemini-ს B დავალებას უნდა დაემატოს (რეკომენდაცია A22).
  \item \textbf{ფართობით შეწონილი გასაშუალოება.} ესაა \textbf{მეორე სტატია} (Suttirat-ის შემდეგ), რომელიც გრიდულ კლიმატს
    რეგიონის მთელ ფართობზე აშუალოებს. შენი სისტემა კი \textbf{ერთ წერტილს} იღებს (ქვეყნის ცენტრი). ბრაზილიისნაირ
    დიდ ქვეყანაზე (8,5 მლნ კმ$^2$, ტროპიკები + ზომიერი სამხრეთი) ერთი წერტილი მეცნიერულად გამოუსადეგარია. ეს რეკომენდაცია A6-ს
    (მრავალწერტილოვანი კლიმატი) აძლიერებს; თეზისში შეზღუდვად აუცილებლად ახსენე.
  \item \textbf{Tmin/Tmax}: მეორე სტატია, რომელიც Tmin-ს იყენებს და აღმოაჩენს, რომ Tmin Tmax-ზე უკეთესია (გვ.~e215).
    დავალება C (Tmin/Tmax) ამით კიდევ უფრო დასაბუთებულია.
\end{enumerate}
\end{alarm}

\begin{termin}
\textbf{CRU TS (Climatic Research Unit gridded Time Series)}: აღმოსავლეთ ინგლისის უნივერსიტეტის გრიდული კლიმატის მონაცემები,
სადგურების დაკვირვებებიდან ინტერპოლირებული (არა რეანალიზი, როგორც ERA5).\\
\textbf{PDSI / scPDSI}: გვალვის ინდექსი; sc = self-calibrated, თვითკალიბრებადი.\\
\textbf{exactextractr}: R-ის პაკეტი, რომელიც გრიდს პოლიგონზე (რეგიონზე) ზუსტად, ფართობით შეწონვით აშუალოებს.
\end{termin}

\subsection*{ურბანიზაცია და წყალზე წვდომა} \pg{e211--e212}
იმის შესაფასებლად, განსხვავდება თუ არა კავშირი ურბანიზაციისა და წყალმომარაგების მიხედვით, DATASUS-იდან ავიღეთ
\textbf{2010 წლის აღწერის} მონაცემები: ქალაქად მცხოვრებთა წილი და წყალსადენზე წვდომის მქონე მცხოვრებთა წილი. ცუდი სანიტარია
კოღოების გამრავლებას უწყობს ხელს, ამიტომ მოსალოდნელია, რომ ასეთ რაიონებში რისკი მაღალია. თუმცა მიკრორეგიონის დონეზე
ქალაქად მცხოვრებთა წილი \textbf{დადებითად კორელირებს} წყალსადენზე წვდომასთან, ამიტომ ამ დონის აგრეგაციისას წყალზე წვდომის
ცვლადი არ გამოდგება \textbf{კოლინეარობის} გამო (\textbf{პირსონის კორელაციის კოეფიციენტი r = 0,73}, $p < 0{,}0001$).

წყალმომარაგების ხარისხი და საიმედოობა ძნელად იზომება. ერთ-ერთი გზაა სისტემის შეფერხებების მონიტორინგი, რომელსაც მომწოდებლები
ყოველწლიურად აცხადებენ \textbf{სანიტარიის ეროვნულ საინფორმაციო სისტემაში} (SNIS). ამ კვლევისთვის 2000--2016 წლებში
მუნიციპალიტეტების მიერ დაფიქსირებული შეფერხებების რაოდენობა \textbf{გაიყო წლებისა და მუნიციპალიტეტების რაოდენობაზე} თითოეულ
მიკრორეგიონში, რითაც მივიღეთ შეფერხებების სიხშირე (0-დან 1-მდე). ეს ცვლადი ურბანიზაციასთან სუსტად კორელირებს
(r = 0,13, $p = 0{,}0017$).

\begin{etl}
\textbf{ორი საინტერესო რამ შენთვის:}
\begin{enumerate}[leftmargin=1.5em]
  \item \textbf{ავტორები პირსონის კორელაციას ცვლადების შესარჩევად იყენებენ} (r = 0,73 $\rightarrow$ ცვლადი ამოიღეს). შენი
    სისტემის კორელაციის ფუნქცია (FR-19) ზუსტად ამ ტიპის \textbf{წინასწარ, საორიენტაციო ანალიზს} ემსახურება. ეს კარგი არგუმენტია:
    ის მოდელის ნაცვლად კი არა, \textbf{მოდელამდე} გამოიყენება.
  \item \textbf{სხვადასხვა დროითი რეზოლუცია ერთ მოდელში}: თვიური (შემთხვევები, კლიმატი), წლიური (მოსახლეობა), სტატიკური
    (2010 წლის აღწერა) და 17-წლიანი საშუალო (წყლის შეფერხებები). ამ ყველაფრის ერთ ცხრილში მოყვანა ზუსტად ჰარმონიზაციის
    ამოცანაა. შენი სისტემა ამას ნაწილობრივ აკეთებს (თვიური + წლიური); სტატიკური სოციალურ-ეკონომიკური ცვლადები ჯერ არ აქვს
    (Future Work, A9/A11-ის მსგავსად).
\end{enumerate}
\end{etl}

\subsection*{მოდელირების მიდგომა \textit{[შემოკლებით]}} \pg{e212--e214}
\begin{skipbox}
ეს ნაწილი სტატისტიკური მოდელის აღწერაა; შენი სისტემა მოდელს \textbf{არ} აკეთებს, მხოლოდ მონაცემებს ამზადებს მისთვის.
საკმარისია იცოდე:
\begin{itemize}[leftmargin=1.5em]
  \item მოდელი: სივრცულ-დროითი \textbf{ბაიესური იერარქიული მოდელი}, \textbf{უარყოფითი ბინომიალური} განაწილებით (შემთხვევების
    ჭარბი ვარიაციის გამო), შეფასება R-INLA-თი.
  \item შემთხვევითი ეფექტები: შტატის დონის თვიური (სეზონურობა) და მიკრორეგიონის დონის წლიური (წლებს შორის ცვალებადობა).
  \item DLNM Tmin-ისთვის, Tmax-ისთვის და PDSI-სთვის, \textbf{lag 0--6 თვე}.
  \item PDSI-სა და ურბანიზაციას შორის წრფივი ურთიერთქმედება; ურბანიზაცია სამ დონეზე ცენტრირებული (58\%, 73\%, 87\%).
  \item მოდელის შერჩევა: DIC და ჯვარედინი ვალიდაციის log score; ჯვარედინი ვალიდაცია 19 $\times$ 12-ჯერ (ყოველ ჯერზე წლის ერთი თვე
    გამორიცხული).
\end{itemize}
\end{skipbox}

\begin{citat}
„DLNMs were used to explore possible non-linear and delayed associations between dengue incidence, temperature (Tmin and Tmax), and
the PDSI from 0 to 6 months.“ \pg{e212}
\end{citat}

%==========================================================
\section{Results (შედეგები) \textit{[შემოკლებით]}} \pg{e214--e215}

მონაცემთა ნაკრები 558 მიკრორეგიონში 2001--2019 წლებში 12\,895\,293 შემთხვევას მოიცავდა. 19 წლის განმავლობაში ინციდენტობა
გაიზარდა და გადაცემის ზონა სამხრეთით, ცენტრალურ-დასავლეთ რეგიონსა და ამაზონიაში გაფართოვდა. სეზონურობა ქვეყნის მასშტაბით
განსხვავდება: ჩრდილოეთში პიკი წლის დასაწყისშია, ჩრდილო-აღმოსავლეთში გვიან. დიდი ეპიდემიები იყო 2010, 2013, 2015, 2016 და
2019 წლებში.

Tmin-ისა და PDSI-ს ჩართვამ მოდელი ყველაზე მეტად გააუმჯობესა; \textbf{Tmin უკეთესი აღმოჩნდა, ვიდრე Tmax} (გვ.~e215).
რისკი Tmin-ის ზრდასთან ერთად თანდათან იზრდებოდა; მაქსიმუმი Tmin = 25,5$^\circ$C-ზე, 2--4 თვის lag-ით. შერჩეულმა მოდელმა
558-დან 409 (73\%) მიკრორეგიონში საბაზისო მოდელზე უკეთესი შედეგი აჩვენა. მოდელმა სწორად ამოიცნო 2010, 2013, 2015--16 და 2019
წლების ეპიდემიები და დაბალი ინციდენტობის წლები (2004, 2017).

\begin{center}
\begin{tabular}{@{}lcccccc@{}}
\toprule
 & \multicolumn{3}{c}{\textbf{ნოტიო პირობები}} & \multicolumn{3}{c}{\textbf{მშრალი პირობები}}\\
 & PDSI & lag (თვე) & RR & PDSI & lag (თვე) & RR\\
\midrule
საერთო & 7,0 & 1 & 1,56 & --7,5 & 4 & 1,43\\
მაღალი ურბანიზაცია (87\%) & 7,0 & 2 & 1,58 & --7,5 & 4 & 1,60\\
საშუალო (73\%) & 7,0 & 1 & 1,56 & --7,5 & 4 & 1,40\\
დაბალი (58\%) & 7,0 & 0 & 1,77 & --4,0 & 4 & 1,15\\
\bottomrule
\end{tabular}\\[0.3em]
{\small სტატიის Table (გვ.~e216), 95\% CI-ის გარეშე.}
\end{center}

\begin{etl}
\textbf{lag-ების დიაპაზონი:} Suttirat 1--3 თვეს იყენებს, Lowe \textbf{0--6} თვეს; გვალვის ეფექტი \textbf{4 თვის} შემდეგაა
მაქსიმალური. შენი FR-19 ამჟამად მხოლოდ lag 0--3-ს ითვლის. ამ სტატიის მიხედვით 0--3 \imp{ზოგ შემთხვევაში არასაკმარისია}.
Gemini-ს B დავალება lag-ებს 1--6 დიაპაზონში უშვებს, რაც ამ სტატიით დასაბუთებულია. კორელაციის lag-ების (FR-19)
6-მდე გაფართოებაც შეიძლება განიხილო.
\end{etl}

%==========================================================
\section{Discussion (დისკუსია) \textit{[შემოკლებით]}} \pg{e215--e217}

ავტორები ხსნიან, რატომ ჩნდება დაგვიანება: გვალვის შემდეგ ხალხი თანდათან იწყებს წყლის შენახვას იმპროვიზირებულ ჭურჭელში;
\textit{Aedes aegypti}-ს კვერცხები მშრალ პირობებში 120 დღე ძლებს. ძლიერი წვიმის შემდეგ ლარვების საბინადრო მატულობს და
რამდენიმე კვირაში კოღოების პოპულაცია იზრდება. ავტორები ხაზს უსვამენ, რომ მოდელები \textbf{ღია წვდომის მონაცემთა ბაზით}
(12 მილიონზე მეტი შემთხვევა) შეიქმნა, და გვალვის მონიტორინგისა და სეზონური კლიმატის პროგნოზების ადრეული გაფრთხილების
სისტემებში ჩართვას გვთავაზობენ.

\begin{citat}
„This study provides additional evidence to motivate the incorporation of drought monitoring and seasonal climate forecasts in
vector surveillance and dengue early warning systems [\ldots]“ \pg{e216}
\end{citat}

\subsection*{შეზღუდვები: სრული თარგმანი} \pg{e217}
ეროვნული მონაცემების გამოყენებას უპირატესობა აქვს კლიმატური და სოციალურ-ეკონომიკური პირობების ფართო დიაპაზონის შესასწავლად,
მაგრამ ნაკლიც აქვს \textbf{მონაცემთა ხარისხისა} და დენგეს ნამდვილი ტვირთის ასახვის თვალსაზრისით. შემთხვევები SINAN-იდანაა,
რომელიც \textbf{პასიური ზედამხედველობის} სისტემაა, ამიტომ მსუბუქი ან უსიმპტომო ინფექციები, რომლებიც, სავარაუდოდ, შემთხვევების
უმეტესობაა, შეიძლება გამოგვრჩეს. ამასთან, დარეგისტრირებული შემთხვევების მხოლოდ მცირე ნაწილია \textbf{ლაბორატორიულად
დადასტურებული} (10\%-დან ჩრდილო-აღმოსავლეთში 30\%-მდე სამხრეთში), ეპიდემიების დროს კი შესაძლოა უფრო ნაკლები. ერთი კვლევის
შეფასებით, დაბალი გადაცემის პერიოდში 40 შემთხვევიდან მხოლოდ ერთი ფიქსირდებოდა. ლაბორატორიული დადასტურების არარსებობა
\textbf{არასწორი კლასიფიკაციის} რისკს ზრდის, განსაკუთრებით 2016 წლიდან, ზიკასა და ჩიკუნგუნიას გავრცელების შემდეგ (იგივე
გადამტანი, მსგავსი სიმპტომები). შემთხვევები და მოსახლეობა ასაკისა და სეროტიპის მიხედვით არ არის დაყოფილი.

ურბანიზაციის ცვლადი \textbf{სტატიკურია} (2010 წლის აღწერა) და ბოლო ათწლეულში, სავარაუდოდ, შეიცვალა. წყლის შეფერხებების
ცვლადი დამოკიდებულია იმაზე, აცხადებენ თუ არა მომწოდებლები შეფერხებებს, და ანგარიშგების შეცდომებისადმი მგრძნობიარეა.
ქალაქებში შესაძლოა შემთხვევები უკეთ ფიქსირდებოდეს (სამედიცინო დაწესებულებებზე უკეთესი წვდომა). კოღოების სიმჭიდროვის მონაცემები
მიკრორეგიონის დონეზე ხელმისაწვდომი არ იყო.

\begin{alarm}
\textbf{შემთხვევის განსაზღვრება და მონაცემთა ხარისხი.} ეს უკვე მესამე წყაროა (Clarke 2024-ისა და Suttirat-ის შემდეგ), რომელიც
აჩვენებს, რომ „შემთხვევა“ სხვადასხვა რამეს ნიშნავს: დადასტურებული 10--30\%, დანარჩენი კლინიკური დიაგნოზია; 2016 წლიდან ზიკა
და ჩიკუნგუნია ერევა. ეს რეკომენდაცია \textbf{A15ა}-ს (გაფრთხილება PDF-სა და UI-ში) აძლიერებს. შეგიძლია ორივე წყარო ციტირო.
\end{alarm}

%==========================================================
\section{Data sharing (მონაცემთა გაზიარება): სრული თარგმანი} \pg{e218}

ამ კვლევაში გამოყენებული ყველა მონაცემი ღია წვდომისაა და ინტერნეტში თავისუფლად ხელმისაწვდომია (დეტალები მეთოდებში).
ანალიზისთვის გამოყენებული მონაცემები და კოდი ხელმისაწვდომია
\href{https://github.com/drrachellowe/hydromet_dengue}{github.com/drrachellowe/hydromet\_dengue} და დაარქივებულია მუდმივ
რეპოზიტორიუმში (Zenodo, DOI 10.5281/zenodo.4632205).

\textbf{ავტორების წვლილი (ამონარიდი):} „SAL collated and managed the database and helped with visualisation and drafting of the
manuscript. CB, RdCC, and MSC collected data and did a literature search.“ \pg{e218}

\begin{citat}
„All data used in this study is open access and freely available on the internet, see the methods section for details.“ \pg{e218}
\end{citat}

\begin{etl}
\textbf{ეს მოტივაციისთვის ძალიან მნიშვნელოვანია, მაგრამ ფრთხილად უნდა ჩამოაყალიბო.}
\begin{itemize}[leftmargin=1.5em]
  \item \textbf{ფაქტი:} ყველა მონაცემი ღია იყო, მაგრამ \textbf{ხუთი სხვადასხვა წყაროდან} (DATASUS/SINAN, IBGE, CRU TS, 2010 წლის
    აღწერა, SNIS), სხვადასხვა ფორმატით (ცხრილები, გრიდი) და რეზოლუციით (თვიური, წლიური, სტატიკური).
  \item \textbf{ფაქტი:} ავტორთა წვლილის აღწერაში მონაცემთა ბაზის შეგროვება და მართვა \textbf{ცალკე ამოცანადაა} დასახელებული
    (ერთი ავტორი: „collated and managed the database“; სამი ავტორი: „collected data“).
  \item \textbf{შენი ინტერპრეტაცია (არა ავტორების სიტყვები!):} მონაცემების „ღიაობა“ არ ნიშნავს, რომ ისინი \textbf{გამოსაყენებლად
    მზადაა}. მათი პოვნა, ჩამოტვირთვა, აგრეგაცია და გაერთიანება მაინც საჭიროა და ცალკე ადამიანის სამუშაოა. სწორედ ამ ნაბიჯს
    ამარტივებს შენი სისტემა. თეზისში ასე დაწერე: \textit{„Obwohl alle Daten offen verfügbar waren, stammten sie aus fünf
    verschiedenen Quellen [\ldots]; die Datenaufbereitung wird in den Autorenbeiträgen als eigene Aufgabe ausgewiesen (Lowe et al.
    2021, S.~e218).“} ავტორებს არ მიაწერო, რომ ეს „რთული“ ან „შრომატევადი“ იყო: ეს მათ არ დაუწერიათ.
  \item \textbf{GitHub-ის რეპოზიტორიუმი:} მონაცემები და R კოდი ღიაა. შეგიძლია ნახო, როგორ მოამზადეს მონაცემები, და Case Study-სთვის
    გამოიყენო (მაგ. ბრაზილიის ეროვნულ დონეზე შეკრებილი შემთხვევები შენი სისტემის OpenDengue-ს მონაცემებთან შეადარო).
\end{itemize}
\end{etl}

%==========================================================
\section{შედარება: სტატიის ნაბიჯები vs. შენი სისტემა}\label{sec:shedareba}

\begin{longtable}{@{}p{0.04\linewidth}p{0.36\linewidth}p{0.38\linewidth}p{0.12\linewidth}@{}}
\toprule
\textbf{\#} & \textbf{ნაბიჯი სტატიაში} & \textbf{შენი სისტემა} & \textbf{სტატუსი}\\
\midrule
\endhead
1 & თვიური შემთხვევები, DATASUS/SINAN, 558 მიკრორეგიონი & OpenDengue, მხოლოდ ქვეყნის დონე & ნაწილობრივ\\
2 & წლიური მოსახლეობა (IBGE) & World Bank / UN WPP / ატვირთვა & $\checkmark$\\
3 & ინციდენტობა 100\,000-ზე თვეში & FR-17 & $\checkmark$\\
4 & გრიდული კლიმატი (CRU TS 4.04, 0,5$^\circ$) & ERA5 / NASA POWER, ერთი წერტილი & ნაწილობრივ\\
5 & Tmin, Tmax (დღიური მინ./მაქს. თვიური საშუალო) & მხოლოდ საშუალო ტემპ. (დავალება C) & $\times$ (დაგეგმილი)\\
6 & ფართობით შეწონილი აგრეგაცია რეგიონზე (exactextractr) & ერთი წერტილი & $\times$ (A6)\\
7 & კლიმატი 1 წლით ადრე, lag-ისთვის & არა & $\times$ (A22, ახალი)\\
8 & გვალვის ინდექსი (scPDSI) & არა & $\times$ (A23, Future Work)\\
9 & lag 0--6 თვე & კორელაცია lag 0--3; lag სვეტები 1--6 (დავალება B) & ნაწილობრივ\\
10 & სტატიკური სოციალურ-ეკონომიკური ცვლადები (2010 აღწერა) & არა & $\times$ (Future Work)\\
11 & წყლის შეფერხებები: ჯამი / (წლები $\times$ მუნიციპალიტეტები) & არა & $\times$ (Future Work)\\
12 & კოლინეარობის შემოწმება პირსონის r-ით & Pearson + Spearman (FR-19) & $\checkmark$ (მსგავსი)\\
13 & მონაცემები + კოდი GitHub/Zenodo-ზე & provenance / რეპროდუცირებადობა (A5, A16) & ნაწილობრივ\\
14 & ბაიესური მოდელი, DLNM, INLA & სისტემის ფარგლებს გარეთაა (მოდელს არ აკეთებს) & ---\\
\bottomrule
\end{longtable}

\begin{etl}
\textbf{შეჯამება არგუმენტისთვის:} ნაბიჯები 1--7, 9 და 12 \textbf{ორივე} სტატიაში გვხვდება (Suttirat 2026 და Lowe 2021):
შემთხვევები + მოსახლეობა $\rightarrow$ ინციდენტობა, გრიდული კლიმატი $\rightarrow$ რეგიონზე აგრეგაცია, Tmin/Tmax, lag-ები,
კორელაციით ცვლადების შერჩევა. ეს თეზისში ასე შეგიძლია ჩამოაყალიბო: \textit{„Die Datenaufbereitungsschritte wiederholen sich
studienübergreifend“} და მოთხოვნების კატალოგში ეს ორივე წყარო მიუთითო (traceability).
\end{etl}

%==========================================================
\section{შეჯამება შენთვის}

\subsection*{ციტატები თეზისისთვის}
\begin{longtable}{@{}p{0.6\linewidth}p{0.1\linewidth}p{0.24\linewidth}@{}}
\toprule
\textbf{ციტატა (შემოკლებით)} & \textbf{გვერდი} & \textbf{რას ამტკიცებს}\\
\midrule
„The effect of climate variability\ldots is complex, non-linear, and often delayed by several weeks to months\ldots“ & e211 & lag-ები; Pearson-ის შეზღუდვა\\
„We obtained monthly notified dengue cases\ldots from\ldots DATASUS\ldots yearly population estimates\ldots from\ldots IBGE“ & e211 & ნაბიჯი ტიპურია (FR-17)\\
„\ldots data for 2000 were extracted to allow for a lag period before the first dengue observation\ldots“ & e211 & მოთხოვნა A22\\
„The gridded datasets were aggregated to each microregion\ldots weighted by the percentage\ldots“ & e211 & სივრცული აგრეგაცია (A6)\\
„DLNMs were used to explore\ldots temperature (Tmin and Tmax), and the PDSI from 0 to 6 months.“ & e212 & Tmin/Tmax, lag 0--6\\
„All data used in this study is open access and freely available on the internet\ldots“ & e218 & ღია $\neq$ მზა (ინტერპრეტაცია!)\\
„SAL collated and managed the database\ldots“ & e218 & მონაცემთა მართვა ცალკე ამოცანაა\\
\bottomrule
\end{longtable}

\subsection*{რა შეამოწმო ან დაამატო შენს სისტემაში}
\begin{longtable}{@{}p{0.05\linewidth}p{0.58\linewidth}p{0.08\linewidth}p{0.23\linewidth}@{}}
\toprule
\textbf{\#} & \textbf{საკითხი} & \textbf{პრ.} & \textbf{წყარო სტატიაში}\\
\midrule
1 & \textbf{lag buffer}: lag სვეტებისთვის კლიმატი \texttt{start\_date}-ზე $k$ თვით ადრე ჩამოიტვირთოს, რომ პირველი თვეები ცარიელი არ იყოს (დავალება B-ს დამატება, A22) & P0 & e211\\
2 & Tmin/Tmax (დავალება C) ახლა ორი სტატიითაა დასაბუთებული & P0 & e211, e215\\
3 & კორელაციის lag-ების 0--3-დან 0--6-მდე გაფართოების განხილვა (FR-19) & P2 & e212, e216\\
4 & ერთი წერტილის შეზღუდვა თეზისში: ბრაზილიის მაგალითით (8,5 მლნ კმ$^2$) & P1 & e211\\
5 & გვალვის ინდექსი (მაგ. SPEI/PDSI) როგორც ახალი ცვლადი: Future Work (A23) & P2 & e211\\
6 & შემთხვევის განსაზღვრების გაფრთხილება (A15ა): მესამე წყარო & P1 & e217\\
\bottomrule
\end{longtable}

\subsection*{შესაძლო Case Study}
ბრაზილია შეიძლება მეორე ან მესამე Case Study იყოს: მონაცემები და კოდი ღიაა (GitHub/Zenodo), OpenDengue-ში ბრაზილია არის.
შენი სისტემით ბრაზილიის ეროვნული თვიური ცხრილის შექმნა და ავტორების მონაცემებთან (ეროვნულ დონეზე შეკრებილთან) შედარება
კარგად აჩვენებს, რომელი ნაბიჯები ავტომატიზდება და რომელი არა (მიკრორეგიონები, PDSI). \textbf{ჯერ ხელმძღვანელებთან შეათანხმე}
(რეკომენდაცია B6).

\end{document}
